% The Reformed School — LuaLaTeX, memoir class, 6 x 9 in.
\documentclass[11pt,twoside,openright]{memoir}

% ---------- page geometry: 6 x 9 trim ----------
\setstocksize{9in}{6in}
\settrimmedsize{9in}{6in}{*}
\settrimsize{*}{*}{*}
\settypeblocksize{6.95in}{4.15in}{*}
\setlrmargins{0.95in}{*}{*}     % inner (spine) margin; outer computed
\setulmargins{0.85in}{*}{*}     % upper margin; lower computed
\setheadfoot{\onelineskip}{2\onelineskip}
\setheaderspaces{*}{1.5\onelineskip}{*}
\checkandfixthelayout

% ---------- fonts ----------
\usepackage{fontspec}
\setmainfont{EB Garamond}[
  Numbers={OldStyle,Proportional},
  Ligatures={TeX,Common},
  ItalicFont={EB Garamond Italic},
]
\usepackage{microtype}
\usepackage{ragged2e}
\setlength{\parindent}{1.1em}
\setlength{\parskip}{0pt}
\linespread{1.06}
\usepackage[english]{babel}
\usepackage{csquotes}
\usepackage{hyperref}
\hypersetup{hidelinks,pdftitle={The Reformed School},pdfauthor={},bookmarksnumbered=false}

% ---------- running heads ----------
\makepagestyle{rs}
\makeevenhead{rs}{\small\thepage}{\small\scshape\MakeLowercase{The Reformed School}}{}
\makeoddhead{rs}{}{\small\itshape\rightmark}{\small\thepage}
\makeevenfoot{rs}{}{}{}
\makeoddfoot{rs}{}{}{}
\makeatletter
\makepsmarks{rs}{%
  \createmark{chapter}{right}{nonumber}{}{}%
}
\makeatother
\pagestyle{rs}
\copypagestyle{chapter}{plain}
\makeevenfoot{plain}{}{\small\thepage}{}
\makeoddfoot{plain}{}{\small\thepage}{}

% ---------- chapter (text) openings ----------
\makechapterstyle{rs}{%
  \setlength{\beforechapskip}{2.2\onelineskip}
  \setlength{\midchapskip}{0.6\onelineskip}
  \setlength{\afterchapskip}{2\onelineskip}
  \renewcommand*{\chapnamefont}{\normalfont\small\scshape}
  \renewcommand*{\chapnumfont}{\normalfont\small\scshape}
  \renewcommand*{\printchaptername}{\chapnamefont\MakeLowercase{text}}
  \renewcommand*{\chapternamenum}{\space}
  \renewcommand*{\printchapternum}{\chapnumfont\thechapter}
  \renewcommand*{\chaptitlefont}{\normalfont\Large\itshape}
  \renewcommand*{\printchaptertitle}[1]{\chaptitlefont ##1}
  \renewcommand*{\afterchapternum}{\par\nobreak\vskip\midchapskip}
}
\chapterstyle{rs}
\setsecheadstyle{\normalfont\small\scshape\MakeLowercase}
\setsubsecheadstyle{\normalfont\itshape}
\setbeforesecskip{-1.4\onelineskip}
\setaftersecskip{0.5\onelineskip}
\setsecnumdepth{chapter}
\maxtocdepth{chapter}

% ---------- apparatus environments ----------
% headnote: the editor's introduction to a text, set smaller and ragged
\newenvironment{headnote}{%
  \small\setlength{\parindent}{0pt}\setlength{\parskip}{0.45\onelineskip}%
  \RaggedRight}{%
  \par\vspace{0.6\onelineskip}%
  \noindent\hfill{\normalsize\normalfont$\ast\quad\ast\quad\ast$}\hfill\null\par
  \vspace{0.9\onelineskip}\normalsize}
% source line: edition statement printed under the text's title
\newcommand{\sourceline}[1]{\noindent{\small\itshape #1}\par\vspace{1.1\onelineskip}}
% cut: an omission in an abridged text
\newcommand{\cut}{\par\vspace{0.35\onelineskip}\noindent\hfill$\cdot\quad\cdot\quad\cdot$\hfill\null\par\vspace{0.35\onelineskip}}
% editorial insertion in a text
\newcommand{\ed}[1]{[#1]}
% glossary of words in a headnote
\newenvironment{gloss}{\begin{description}\setlength{\itemsep}{0pt}}{\end{description}}
% a note appended to a text (translator's or editor's notes)
\newenvironment{textnotes}{\par\vspace{\onelineskip}\small\noindent\textsc{notes}\par\nobreak\vspace{0.3\onelineskip}\setlength{\parindent}{0pt}\setlength{\parskip}{0.3\onelineskip}}{\normalsize}
% pandoc compatibility
\providecommand{\tightlist}{\setlength{\itemsep}{0pt}\setlength{\parskip}{0pt}}

\title{The Reformed School}
\begin{document}

% ================= FRONT MATTER =================
\frontmatter
\pagestyle{empty}
\begin{center}
\vspace*{2.8in}
{\Large\itshape The Reformed School}
\end{center}

\cleardoublepage
\begin{center}
\vspace*{1.6in}
{\Huge\itshape The Reformed School\par}
\vspace{0.5in}
{\large Founding Texts of Reformed Education,\\ from Luther to Machen\par}
\vspace{2.2in}
{\normalsize\scshape\MakeLowercase{Privately printed}\par}
{\small \the\year\par}
\end{center}

\clearpage
\vspace*{\fill}
\begin{flushleft}\small
\textit{The Reformed School: Founding Texts of Reformed Education, from Luther to Machen}\\[0.5\onelineskip]
Privately printed. Not for sale.\\[0.5\onelineskip]
The texts reprinted in this book were published between 1524 and 1929 and are in the public domain, as are the English translations used, all of which were published before 1931. Each text is printed from the edition named in its headnote and in the note on sources at the end of the book, and was checked against the page images of that edition. Spelling and punctuation are as printed, except where a headnote says that spelling has been modernised.\\[0.5\onelineskip]
The introduction, headnotes, and afterword are the editor's own.\\[0.5\onelineskip]
Typeset in EB Garamond with Lua\TeX{} and the memoir class. Trim size 6 by 9 inches.
\end{flushleft}

\cleardoublepage
\pagestyle{rs}
\tableofcontents*
\cleardoublepage
\input{front/introduction}

% ================= THE TEXTS =================
\mainmatter
\input{texts.tex}

% ================= BACK MATTER =================
\backmatter
\input{back/afterword}
\input{back/sources}

\end{document}
